\scanpage{12}
% archive foliation: 11
$\Lambda^k(\Lambda^i\Phi_{*k})$ et $\Lambda^k(\Lambda^j\Psi_{*k})$ sont spéciaux
(bien sûr, ils seront même str. spéciaux!).

Leurs composantes sont évidemment
libres. De plus, comme $\Lambda^i\Phi_{*k}$ est
homotope à $0$, de même $\Lambda^k(\Lambda^i\Phi_{*k})$, et
leurs $\pi_n$ sont donc nuls et a fortiori
libres de \unsure{t.f.|r.f.}.
\margin{?}
]

\uline{Remarque}. Soit $L_*$ $k$-module $\frac12$ simplicial
\deleted{\unsure{str.}} spécial, alors $L_*$ est connu à isom. près,
i.e. les $\mu_i=\operatorname{rg}M_i$ $(i\geq1)$ et $\nu_j=\operatorname{rg}N_j$
\deleted{$j>0$} \inserted{$j\geq0$}
connus, quand on connaît la \deleted{\unsure{suite des}}
$n\longmapsto\operatorname{rg}_kL_n=f_{L_*}(n)$,
\deleted{En effet} et les
$\operatorname{rg}_k\pi_j(L_*)$ $(=\operatorname{rg}_kN_j=\nu_j)$.
En effet,
on aura
\[
\sum_{i\geq1}\mu_i\underset{\unsure{P_i(n+1)}}{\underbrace{\binom{n+1}{i}}}
=f_{L_*}(n)-\sum_{j\geq0}\nu_j\binom{n}{j}
\qquad(n\text{ grand})
\]
\deleted{On en sait}
i.e.
\[
f_{L_*}(n-1)-\sum_{j\geq0}\unsure{\nu_j|\beta_j}P_j(n-1)
=\sum_{i\geq1}\mu_iP_i(n)
\]
on sait que les \deleted{\unsure{$\binom{n}{i}$}} \inserted{$P_i(n)$}
sont des
\unsure{lin. indépendants} et $=$ forment une
base du $\mathbf Z$-module des polynômes
à \deleted{valeurs} \inserted{coeff. entiers} entières sur les entiers.
Donc la connaissance de $\sum_{i\geq1}\mu_iP_i$
pour $n$ grand implique
